%% =====================================================================
%%  T-12-05  The Courtier
%%  -------------------------------------------------------------------
%%  One idea per file. Core is mandatory. Slots in canonical order.
%%  Max 700 words. No layout commands. No filler. New source material
%%  for this entry goes in sources/<BOOK>/T-12-05.tex -- never by
%%  rewriting this file (docs/RULES.md R0.1).
%%  =====================================================================
%% @kind: topic
%% @id: T-12-05
%% @title: The Courtier
%% @subtitle: Living inside somebody else's hierarchy without becoming furniture
%% @chapter: c12
%% @order: 50
%% @status: stable
%% @sources: B3
%% @updated: 2026-09-19
%% @allow-rewrite: B3 ingest 2026-09-19 -- committed scaffold replaced by the written entry (planned -> stable lifecycle)

\topic{T-12-05}{The Courtier}{Living inside somebody else's hierarchy without becoming furniture}

\begin{core}
A court is any hierarchy where one person's moods set everyone's weather. Surviving and advancing inside one is a craft with rules: master your own face before anything else, move at the court's rhythm, never outshine the sovereign, and make yourself pleasant to serve. The courtier is not a flatterer; the courtier is a professional at being usefully agreeable without selling judgement.
\end{core}
\begin{mechanism}
Courts run on proximity and mood, so the first skill is emotional self-command: cry and laugh on command, disguise frustration, fake contentment --- be the master of your own face, because every leak is ammunition for rivals. The second is fitness to the times: a slight affectation of a past era may charm, but last decade's fashion is ludicrous; spirit and thinking must keep up with the court you are actually in, not the one you wish you were in. The mechanism underneath is information: a courtier who neither leaks nor bristles collects everything --- who is rising, what the sovereign fears --- while spending nothing. Agreeableness is the interface; the value is the intelligence and the access it keeps open.
\end{mechanism}
\begin{conditions}
Strongest wherever one principal controls your fate --- royal courts, founder-led firms, small boats with one captain. The craft degrades in flat organisations and partner structures, where alliances (T-12-01) outperform courtiership. Fatal where the court demands output: a courtier with nothing to deliver is decoration, and decoration is first overboard.
\end{conditions}
\begin{application}
  \item Learn the sovereign's actual fears and wants --- not the announced ones --- and align your visible work with them.
  \item Rehearse your face. Reactions are decisions, made in advance, so the moment finds them ready.
  \item Match the court's tempo: never ahead of the fashion by more than taste, never behind it by a season.
  \item Give credit loudly, correct privately, and disagree only where the sovereign can afford to be seen losing.
  \item Be pleasant to serve --- easy to ask, quick, reliable. Access is granted to people who are cheap to use.
\end{application}
\begin{feedback}
  \item Landing: the sovereign improvises with you --- asks your view unscripted. Trust is live.
  \item Landing: rivals stop testing you; you have stopped being cheap to attack.
  \item Failing: your candour is quoted across the court by evening. The interface leaked.
  \item Failing: invitations thin out. Proximity is being reassigned.
\end{feedback}
\begin{failure}
The craft's cost is the mask: worn long enough, the performance edits the performer, and a person who has faked contentment for a decade cannot always locate the real kind. Courtiers are also indexed to their sovereign --- exile by boss change is the occupational disease. And agreeableness carries its own sentence: the courtier absorbs the court's judgement as their own, which is precisely the installed consistency the defence chapters warn about (T-24-02).
\end{failure}
\begin{countermeasures}
  \item In your own organisation, watch who never disagrees with you. You have built a court, and courts are where your information goes to die.
  \item Price the mask: if you cannot remember your actual opinion on something that matters, you have been courtier too long.
  \item When a colleague is relentlessly agreeable, promote them for their disagreements or not at all --- otherwise you are training your own echo.
  \item Keep a private record of what you believed before you arrived. It is the only durable defence against a court's slow rewrite.
\end{countermeasures}

\sources{B3}
\seesources
\seealso{T-12-01, T-09-02, T-02-01}
